
\documentclass{article}
\usepackage{colortbl}
\usepackage{makecell}
\usepackage{multirow}
\usepackage{supertabular}

\begin{document}

\newcounter{utterance}

\centering \large Interaction Transcript for game `matchit\_ascii', experiment `similar\_grid\_2', episode 4 with nemotron{-}3{-}ultra.
\vspace{24pt}

{ \footnotesize  \setcounter{utterance}{1}
\setlength{\tabcolsep}{0pt}
\begin{supertabular}{c@{$\;$}|p{.15\linewidth}@{}p{.15\linewidth}p{.15\linewidth}p{.15\linewidth}p{.15\linewidth}p{.15\linewidth}}
    \# & $\;$A & \multicolumn{4}{c}{Game Master} & $\;\:$B\\
    \hline

    \theutterance \stepcounter{utterance}  
    & & \multicolumn{4}{p{0.6\linewidth}}{
        \cellcolor[rgb]{0.9,0.9,0.9}{
            \makecell[{{p{\linewidth}}}]{
                \texttt{\tiny{[P1$\langle$GM]}}
                \texttt{Du deltar i ett samarbetsbaserat gissningsspel.} \\
\texttt{Målet är att ta reda på om detta rutnät och en annan bild som bara jag kan se är samma.} \\
\texttt{Beskriv ditt rutnät först. Sedan kommer jag att ge min beskrivning och vi kan ställa frågor till varandra om rutnäten för att ta reda på om de är samma. Börja nu din korta rutnätsbeskrivning med "BESKRIVNING:" följt av beskrivningen. Lägg inte till något annat. Det här är ditt rutnät:} \\
\texttt{X X X X X} \\
\texttt{X ▢ ▢ ▢ X} \\
\texttt{X ▢ ▢ ▢ X} \\
\texttt{X ▢ ▢ ▢ X} \\
\texttt{X X X X X} \\
            }
        }
    }
    & \\ \\

    \theutterance \stepcounter{utterance}  
    & \multicolumn{4}{p{0.6\linewidth}}{
        \cellcolor[rgb]{0.9,0.9,0.9}{
            \makecell[{{p{\linewidth}}}]{
                \texttt{\tiny{[P1$\rangle$GM]}}
                \texttt{BESKRIVNING: En 5x5{-}ruta med en kant av X runt en 3x3{-}inre yta av tomma rutor (▢).} \\
            }
        }
    }
    & & \\ \\

    \theutterance \stepcounter{utterance}  
    & & & \multicolumn{2}{p{0.3\linewidth}}{
        \cellcolor[rgb]{0.9,0.9,0.9}{
            \makecell[{{p{\linewidth}}}]{
                \texttt{\tiny{[GM$|$GM]}}
                \texttt{continue} \\
            }
        }
    }
    & & \\ \\

    \theutterance \stepcounter{utterance}  
    & & \multicolumn{4}{p{0.6\linewidth}}{
        \cellcolor[rgb]{0.9,0.9,0.9}{
            \makecell[{{p{\linewidth}}}]{
                \texttt{\tiny{[GM$\rangle$P2]}}
                \texttt{Du deltar i ett samarbetsbaserat gissningsspel.} \\
\texttt{Målet är att ta reda på om detta rutnät och en annan bild som bara jag kan se är samma.} \\
\texttt{Beskriv ditt rutnät först. Sedan kommer jag att ge min beskrivning och vi kan ställa frågor till varandra om rutnäten för att ta reda på om de är samma. Börja nu din korta rutnätsbeskrivning med "BESKRIVNING:" följt av beskrivningen. Lägg inte till något annat. Det här är ditt rutnät:} \\
\texttt{X ▢ X X X} \\
\texttt{X ▢ ▢ ▢ X} \\
\texttt{X ▢ ▢ ▢ X} \\
\texttt{X ▢ ▢ ▢ X} \\
\texttt{X ▢ X X X} \\
            }
        }
    }
    & \\ \\

    \theutterance \stepcounter{utterance}  
    & & & \multicolumn{4}{p{0.6\linewidth}}{
        \cellcolor[rgb]{0.9,0.9,0.9}{
            \makecell[{{p{\linewidth}}}]{
                \texttt{\tiny{[GM$\langle$P2]}}
                \texttt{BESKRIVNING: 5x5{-}rutnät med X på periferin och i hörnen. Inre 3x3{-}område (rader 2{-}4, kolumner 2{-}4) är tomt utom (1,1) och (1,3) som är X (nollindexerad). Worten: en ram av X med två X inuti på toppen av den inre ytan.} \\
            }
        }
    }
     \\ \\

    \theutterance \stepcounter{utterance}  
    & & & \multicolumn{2}{p{0.3\linewidth}}{
        \cellcolor[rgb]{0.9,0.9,0.9}{
            \makecell[{{p{\linewidth}}}]{
                \texttt{\tiny{[GM$|$GM]}}
                \texttt{continue} \\
            }
        }
    }
    & & \\ \\

    \theutterance \stepcounter{utterance}  
    & & \multicolumn{4}{p{0.6\linewidth}}{
        \cellcolor[rgb]{0.9,0.9,0.9}{
            \makecell[{{p{\linewidth}}}]{
                \texttt{\tiny{[GM$\rangle$P2]}}
                \texttt{Detta är min BESKRIVNINGBESKRIVNING: En 5x5{-}ruta med en kant av X runt en 3x3{-}inre yta av tomma rutor (▢).} \\
\texttt{Ställ nu en fråga för att ta reda på nya aspekter av mitt rutnät som kan skilja sig från ditt rutnät. Börja med "FRÅGA:" och lägg inte till något annat.} \\
            }
        }
    }
    & \\ \\

    \theutterance \stepcounter{utterance}  
    & & & \multicolumn{4}{p{0.6\linewidth}}{
        \cellcolor[rgb]{0.9,0.9,0.9}{
            \makecell[{{p{\linewidth}}}]{
                \texttt{\tiny{[GM$\langle$P2]}}
                \texttt{FRÅGA: Är alla rutormar inuti den 3x3{-}inre ytan tomma (▢), eller finns det några X inuti den arean?} \\
            }
        }
    }
     \\ \\

    \theutterance \stepcounter{utterance}  
    & & & \multicolumn{2}{p{0.3\linewidth}}{
        \cellcolor[rgb]{0.9,0.9,0.9}{
            \makecell[{{p{\linewidth}}}]{
                \texttt{\tiny{[GM$|$GM]}}
                \texttt{continue} \\
            }
        }
    }
    & & \\ \\

    \theutterance \stepcounter{utterance}  
    & & \multicolumn{4}{p{0.6\linewidth}}{
        \cellcolor[rgb]{0.9,0.9,0.9}{
            \makecell[{{p{\linewidth}}}]{
                \texttt{\tiny{[P1$\langle$GM]}}
                \texttt{Detta är min BESKRIVNINGBESKRIVNING: 5x5{-}rutnät med X på periferin och i hörnen. Inre 3x3{-}område (rader 2{-}4, kolumner 2{-}4) är tomt utom (1,1) och (1,3) som är X (nollindexerad). Worten: en ram av X med två X inuti på toppen av den inre ytan.} \\
\texttt{FRÅGA: Är alla rutormar inuti den 3x3{-}inre ytan tomma (▢), eller finns det några X inuti den arean?Börja ditt svar med "SVAR:" och lägg inte till något annat.} \\
            }
        }
    }
    & \\ \\

    \theutterance \stepcounter{utterance}  
    & \multicolumn{4}{p{0.6\linewidth}}{
        \cellcolor[rgb]{0.9,0.9,0.9}{
            \makecell[{{p{\linewidth}}}]{
                \texttt{\tiny{[P1$\rangle$GM]}}
                \texttt{SVAR: Alla nio rutor i den 3x3{-}inre ytan är tomma (▢). Det finns inga X inuti den arean.} \\
            }
        }
    }
    & & \\ \\

    \theutterance \stepcounter{utterance}  
    & & & \multicolumn{2}{p{0.3\linewidth}}{
        \cellcolor[rgb]{0.9,0.9,0.9}{
            \makecell[{{p{\linewidth}}}]{
                \texttt{\tiny{[GM$|$GM]}}
                \texttt{continue} \\
            }
        }
    }
    & & \\ \\

    \theutterance \stepcounter{utterance}  
    & & \multicolumn{4}{p{0.6\linewidth}}{
        \cellcolor[rgb]{0.9,0.9,0.9}{
            \makecell[{{p{\linewidth}}}]{
                \texttt{\tiny{[P1$\langle$GM]}}
                \texttt{Ställ nu en fråga för att ta reda på nya aspekter av mitt rutnät som kan skilja sig från ditt rutnät. Börja med "FRÅGA:" och lägg inte till något annat.} \\
            }
        }
    }
    & \\ \\

    \theutterance \stepcounter{utterance}  
    & \multicolumn{4}{p{0.6\linewidth}}{
        \cellcolor[rgb]{0.9,0.9,0.9}{
            \makecell[{{p{\linewidth}}}]{
                \texttt{\tiny{[P1$\rangle$GM]}}
                \texttt{FRÅGA: Vad menar du exakt med att (1,1) och (1,3) är X "på toppen av den inre ytan" – ligger dessa X *inuti* den 3x3{-}inre ytan (alltså på positionerna rad 2, kolumn 2 och rad 2, kolumn 4 i det totala 5x5{-}rutnätet) eller ligger de *utomför* den inre ytan?} \\
            }
        }
    }
    & & \\ \\

    \theutterance \stepcounter{utterance}  
    & & & \multicolumn{2}{p{0.3\linewidth}}{
        \cellcolor[rgb]{0.9,0.9,0.9}{
            \makecell[{{p{\linewidth}}}]{
                \texttt{\tiny{[GM$|$GM]}}
                \texttt{continue} \\
            }
        }
    }
    & & \\ \\

    \theutterance \stepcounter{utterance}  
    & & \multicolumn{4}{p{0.6\linewidth}}{
        \cellcolor[rgb]{0.9,0.9,0.9}{
            \makecell[{{p{\linewidth}}}]{
                \texttt{\tiny{[GM$\rangle$P2]}}
                \texttt{SVAR: Alla nio rutor i den 3x3{-}inre ytan är tomma (▢). Det finns inga X inuti den arean.} \\
\texttt{FRÅGA: Vad menar du exakt med att (1,1) och (1,3) är X "på toppen av den inre ytan" – ligger dessa X *inuti* den 3x3{-}inre ytan (alltså på positionerna rad 2, kolumn 2 och rad 2, kolumn 4 i det totala 5x5{-}rutnätet) eller ligger de *utomför* den inre ytan?Börja ditt svar med "SVAR:" och lägg inte till något annat.} \\
            }
        }
    }
    & \\ \\

    \theutterance \stepcounter{utterance}  
    & & & \multicolumn{4}{p{0.6\linewidth}}{
        \cellcolor[rgb]{0.9,0.9,0.9}{
            \makecell[{{p{\linewidth}}}]{
                \texttt{\tiny{[GM$\langle$P2]}}
                \texttt{SVAR: De ligger inuti den 3x3{-}inre ytan. Med nollindexering är de på positionerna (1,1) och (1,3), vilket motsvarar rad 2, kolumn 2 och rad 2, kolumn 4 i ett 1{-}indexerat 5x5{-}rutnät. Ditt rutnät har alltså en helt tom 3x3{-}yta, medan mitt har två X på den övre raden i den inre ytan. Rutnäten är olika.} \\
            }
        }
    }
     \\ \\

    \theutterance \stepcounter{utterance}  
    & & & \multicolumn{2}{p{0.3\linewidth}}{
        \cellcolor[rgb]{0.9,0.9,0.9}{
            \makecell[{{p{\linewidth}}}]{
                \texttt{\tiny{[GM$|$GM]}}
                \texttt{continue} \\
            }
        }
    }
    & & \\ \\

    \theutterance \stepcounter{utterance}  
    & & \multicolumn{4}{p{0.6\linewidth}}{
        \cellcolor[rgb]{0.9,0.9,0.9}{
            \makecell[{{p{\linewidth}}}]{
                \texttt{\tiny{[GM$\rangle$P2]}}
                \texttt{Ställ nu en fråga för att ta reda på nya aspekter av mitt rutnät som kan skilja sig från ditt rutnät. Börja med "FRÅGA:" och lägg inte till något annat.} \\
            }
        }
    }
    & \\ \\

    \theutterance \stepcounter{utterance}  
    & & & \multicolumn{4}{p{0.6\linewidth}}{
        \cellcolor[rgb]{0.9,0.9,0.9}{
            \makecell[{{p{\linewidth}}}]{
                \texttt{\tiny{[GM$\langle$P2]}}
                \texttt{FRÅGA: Har ditt rutnät X på alla yttre kanter (första och sista rad samt första och sista kolumn) precis som mitt?} \\
            }
        }
    }
     \\ \\

    \theutterance \stepcounter{utterance}  
    & & & \multicolumn{2}{p{0.3\linewidth}}{
        \cellcolor[rgb]{0.9,0.9,0.9}{
            \makecell[{{p{\linewidth}}}]{
                \texttt{\tiny{[GM$|$GM]}}
                \texttt{continue} \\
            }
        }
    }
    & & \\ \\

    \theutterance \stepcounter{utterance}  
    & & \multicolumn{4}{p{0.6\linewidth}}{
        \cellcolor[rgb]{0.9,0.9,0.9}{
            \makecell[{{p{\linewidth}}}]{
                \texttt{\tiny{[P1$\langle$GM]}}
                \texttt{SVAR: De ligger inuti den 3x3{-}inre ytan. Med nollindexering är de på positionerna (1,1) och (1,3), vilket motsvarar rad 2, kolumn 2 och rad 2, kolumn 4 i ett 1{-}indexerat 5x5{-}rutnät. Ditt rutnät har alltså en helt tom 3x3{-}yta, medan mitt har två X på den övre raden i den inre ytan. Rutnäten är olika.} \\
\texttt{FRÅGA: Har ditt rutnät X på alla yttre kanter (första och sista rad samt första och sista kolumn) precis som mitt?Börja ditt svar med "SVAR:" och lägg inte till något annat.} \\
            }
        }
    }
    & \\ \\

    \theutterance \stepcounter{utterance}  
    & \multicolumn{4}{p{0.6\linewidth}}{
        \cellcolor[rgb]{0.9,0.9,0.9}{
            \makecell[{{p{\linewidth}}}]{
                \texttt{\tiny{[P1$\rangle$GM]}}
                \texttt{SVAR: Ja, mitt rutnät har X på alla yttre kanter (första och sista rad samt första och sista kolumn).} \\
            }
        }
    }
    & & \\ \\

    \theutterance \stepcounter{utterance}  
    & & & \multicolumn{2}{p{0.3\linewidth}}{
        \cellcolor[rgb]{0.9,0.9,0.9}{
            \makecell[{{p{\linewidth}}}]{
                \texttt{\tiny{[GM$|$GM]}}
                \texttt{continue} \\
            }
        }
    }
    & & \\ \\

    \theutterance \stepcounter{utterance}  
    & & \multicolumn{4}{p{0.6\linewidth}}{
        \cellcolor[rgb]{0.9,0.9,0.9}{
            \makecell[{{p{\linewidth}}}]{
                \texttt{\tiny{[P1$\langle$GM]}}
                \texttt{Ställ nu en fråga för att ta reda på nya aspekter av mitt rutnät som kan skilja sig från ditt rutnät. Börja med "FRÅGA:" och lägg inte till något annat.} \\
            }
        }
    }
    & \\ \\

    \theutterance \stepcounter{utterance}  
    & \multicolumn{4}{p{0.6\linewidth}}{
        \cellcolor[rgb]{0.9,0.9,0.9}{
            \makecell[{{p{\linewidth}}}]{
                \texttt{\tiny{[P1$\rangle$GM]}}
                \texttt{FRÅGA: Finns det några fler X inuti den 3x3{-}inre ytan utöver de två du nämnde på den övre raden (positionerna (1,1) och (1,3) nollindexerat)?} \\
            }
        }
    }
    & & \\ \\

    \theutterance \stepcounter{utterance}  
    & & & \multicolumn{2}{p{0.3\linewidth}}{
        \cellcolor[rgb]{0.9,0.9,0.9}{
            \makecell[{{p{\linewidth}}}]{
                \texttt{\tiny{[GM$|$GM]}}
                \texttt{continue} \\
            }
        }
    }
    & & \\ \\

    \theutterance \stepcounter{utterance}  
    & & \multicolumn{4}{p{0.6\linewidth}}{
        \cellcolor[rgb]{0.9,0.9,0.9}{
            \makecell[{{p{\linewidth}}}]{
                \texttt{\tiny{[GM$\rangle$P2]}}
                \texttt{SVAR: Ja, mitt rutnät har X på alla yttre kanter (första och sista rad samt första och sista kolumn).} \\
\texttt{FRÅGA: Finns det några fler X inuti den 3x3{-}inre ytan utöver de två du nämnde på den övre raden (positionerna (1,1) och (1,3) nollindexerat)?Börja ditt svar med "SVAR:" och lägg inte till något annat.} \\
            }
        }
    }
    & \\ \\

    \theutterance \stepcounter{utterance}  
    & & & \multicolumn{4}{p{0.6\linewidth}}{
        \cellcolor[rgb]{0.9,0.9,0.9}{
            \makecell[{{p{\linewidth}}}]{
                \texttt{\tiny{[GM$\langle$P2]}}
                \texttt{SVAR: Nej, det finns inga fler X inuti den 3x3{-}inre ytan. De enda X:en inuti är de två på den övre raden (nollindexerad rad 1, kolumn 1 och kolumn 3). Resten av de inre 9 rutorna är tomma (▢).} \\
            }
        }
    }
     \\ \\

    \theutterance \stepcounter{utterance}  
    & & & \multicolumn{2}{p{0.3\linewidth}}{
        \cellcolor[rgb]{0.9,0.9,0.9}{
            \makecell[{{p{\linewidth}}}]{
                \texttt{\tiny{[GM$|$GM]}}
                \texttt{continue} \\
            }
        }
    }
    & & \\ \\

    \theutterance \stepcounter{utterance}  
    & & \multicolumn{4}{p{0.6\linewidth}}{
        \cellcolor[rgb]{0.9,0.9,0.9}{
            \makecell[{{p{\linewidth}}}]{
                \texttt{\tiny{[GM$\rangle$P2]}}
                \texttt{Ställ nu en fråga för att ta reda på nya aspekter av mitt rutnät som kan skilja sig från ditt rutnät. Börja med "FRÅGA:" och lägg inte till något annat.} \\
            }
        }
    }
    & \\ \\

    \theutterance \stepcounter{utterance}  
    & & & \multicolumn{4}{p{0.6\linewidth}}{
        \cellcolor[rgb]{0.9,0.9,0.9}{
            \makecell[{{p{\linewidth}}}]{
                \texttt{\tiny{[GM$\langle$P2]}}
                \texttt{FRÅGA: Om vi tittar på de fyra hörnen i den 3x3{-}inre ytan (positionerna (1,1), (1,3), (3,1), (3,3) nollindexerat), vilka av dessa är X och vilka är tomma (▢) i ditt rutnät?} \\
            }
        }
    }
     \\ \\

    \theutterance \stepcounter{utterance}  
    & & & \multicolumn{2}{p{0.3\linewidth}}{
        \cellcolor[rgb]{0.9,0.9,0.9}{
            \makecell[{{p{\linewidth}}}]{
                \texttt{\tiny{[GM$|$GM]}}
                \texttt{continue} \\
            }
        }
    }
    & & \\ \\

    \theutterance \stepcounter{utterance}  
    & & \multicolumn{4}{p{0.6\linewidth}}{
        \cellcolor[rgb]{0.9,0.9,0.9}{
            \makecell[{{p{\linewidth}}}]{
                \texttt{\tiny{[P1$\langle$GM]}}
                \texttt{SVAR: Nej, det finns inga fler X inuti den 3x3{-}inre ytan. De enda X:en inuti är de två på den övre raden (nollindexerad rad 1, kolumn 1 och kolumn 3). Resten av de inre 9 rutorna är tomma (▢).} \\
\texttt{FRÅGA: Om vi tittar på de fyra hörnen i den 3x3{-}inre ytan (positionerna (1,1), (1,3), (3,1), (3,3) nollindexerat), vilka av dessa är X och vilka är tomma (▢) i ditt rutnät?Börja ditt svar med "SVAR:" och lägg inte till något annat.} \\
            }
        }
    }
    & \\ \\

    \theutterance \stepcounter{utterance}  
    & \multicolumn{4}{p{0.6\linewidth}}{
        \cellcolor[rgb]{0.9,0.9,0.9}{
            \makecell[{{p{\linewidth}}}]{
                \texttt{\tiny{[P1$\rangle$GM]}}
                \texttt{SVAR: I mitt rutnät är alla fyra hörnen i den 3x3{-}inre ytan tomma (▢). Det finns inga X på någon av positionerna (1,1), (1,3), (3,1), (3,3) nollindexerat.} \\
            }
        }
    }
    & & \\ \\

    \theutterance \stepcounter{utterance}  
    & & & \multicolumn{2}{p{0.3\linewidth}}{
        \cellcolor[rgb]{0.9,0.9,0.9}{
            \makecell[{{p{\linewidth}}}]{
                \texttt{\tiny{[GM$|$GM]}}
                \texttt{continue} \\
            }
        }
    }
    & & \\ \\

    \theutterance \stepcounter{utterance}  
    & & \multicolumn{4}{p{0.6\linewidth}}{
        \cellcolor[rgb]{0.9,0.9,0.9}{
            \makecell[{{p{\linewidth}}}]{
                \texttt{\tiny{[P1$\langle$GM]}}
                \texttt{Ställ nu en fråga för att ta reda på nya aspekter av mitt rutnät som kan skilja sig från ditt rutnät. Börja med "FRÅGA:" och lägg inte till något annat.} \\
            }
        }
    }
    & \\ \\

    \theutterance \stepcounter{utterance}  
    & \multicolumn{4}{p{0.6\linewidth}}{
        \cellcolor[rgb]{0.9,0.9,0.9}{
            \makecell[{{p{\linewidth}}}]{
                \texttt{\tiny{[P1$\rangle$GM]}}
                \texttt{FRÅGA: Hur ser den mellersta raden (nollindexerad rad 2, alltså den tredje raden i det 5x5{-}rutnätet) ut i din inre 3x3{-}yta – är alla tre rutor tomma (▢ ▢ ▢) eller finns det något X där?} \\
            }
        }
    }
    & & \\ \\

    \theutterance \stepcounter{utterance}  
    & & & \multicolumn{2}{p{0.3\linewidth}}{
        \cellcolor[rgb]{0.9,0.9,0.9}{
            \makecell[{{p{\linewidth}}}]{
                \texttt{\tiny{[GM$|$GM]}}
                \texttt{continue} \\
            }
        }
    }
    & & \\ \\

    \theutterance \stepcounter{utterance}  
    & & \multicolumn{4}{p{0.6\linewidth}}{
        \cellcolor[rgb]{0.9,0.9,0.9}{
            \makecell[{{p{\linewidth}}}]{
                \texttt{\tiny{[GM$\rangle$P2]}}
                \texttt{SVAR: I mitt rutnät är alla fyra hörnen i den 3x3{-}inre ytan tomma (▢). Det finns inga X på någon av positionerna (1,1), (1,3), (3,1), (3,3) nollindexerat.} \\
\texttt{FRÅGA: Hur ser den mellersta raden (nollindexerad rad 2, alltså den tredje raden i det 5x5{-}rutnätet) ut i din inre 3x3{-}yta – är alla tre rutor tomma (▢ ▢ ▢) eller finns det något X där?Börja ditt svar med "SVAR:" och lägg inte till något annat.} \\
            }
        }
    }
    & \\ \\

    \theutterance \stepcounter{utterance}  
    & & & \multicolumn{4}{p{0.6\linewidth}}{
        \cellcolor[rgb]{0.9,0.9,0.9}{
            \makecell[{{p{\linewidth}}}]{
                \texttt{\tiny{[GM$\langle$P2]}}
                \texttt{SVAR: I mitt rutnät är alla tre rutor på den mellersta raden i den inre 3x3{-}ytan (nollindexerad rad 2, positioner (2,1), (2,2), (2,3)) tomma (▢ ▢ ▢). Det finns inga X på den raden.} \\
            }
        }
    }
     \\ \\

    \theutterance \stepcounter{utterance}  
    & & & \multicolumn{2}{p{0.3\linewidth}}{
        \cellcolor[rgb]{0.9,0.9,0.9}{
            \makecell[{{p{\linewidth}}}]{
                \texttt{\tiny{[GM$|$GM]}}
                \texttt{continue} \\
            }
        }
    }
    & & \\ \\

    \theutterance \stepcounter{utterance}  
    & & \multicolumn{4}{p{0.6\linewidth}}{
        \cellcolor[rgb]{0.9,0.9,0.9}{
            \makecell[{{p{\linewidth}}}]{
                \texttt{\tiny{[GM$\rangle$P2]}}
                \texttt{Fatta nu ett beslut. Vad tror du: är ditt rutnät och det andra rutnätet som beskrivs samma? Skriv "BESLUT: samma rutnät" om du tror att de är samma rutnät eller "BESLUT: olika rutnät" om du tror att de är olika rutnät. Lägg inte till något annat.} \\
            }
        }
    }
    & \\ \\

    \theutterance \stepcounter{utterance}  
    & & & \multicolumn{4}{p{0.6\linewidth}}{
        \cellcolor[rgb]{0.9,0.9,0.9}{
            \makecell[{{p{\linewidth}}}]{
                \texttt{\tiny{[GM$\langle$P2]}}
                \texttt{BESLUT: olika rutnät} \\
            }
        }
    }
     \\ \\

    \theutterance \stepcounter{utterance}  
    & & & \multicolumn{2}{p{0.3\linewidth}}{
        \cellcolor[rgb]{0.9,0.9,0.9}{
            \makecell[{{p{\linewidth}}}]{
                \texttt{\tiny{[GM$|$GM]}}
                \texttt{continue} \\
            }
        }
    }
    & & \\ \\

    \theutterance \stepcounter{utterance}  
    & & & \multicolumn{2}{p{0.3\linewidth}}{
        \cellcolor[rgb]{0.9,0.9,0.9}{
            \makecell[{{p{\linewidth}}}]{
                \texttt{\tiny{[GM$|$GM]}}
                \texttt{success} \\
            }
        }
    }
    & & \\ \\

    \theutterance \stepcounter{utterance}  
    & & \multicolumn{4}{p{0.6\linewidth}}{
        \cellcolor[rgb]{0.9,0.9,0.9}{
            \makecell[{{p{\linewidth}}}]{
                \texttt{\tiny{[P1$\langle$GM]}}
                \texttt{SVAR: I mitt rutnät är alla tre rutor på den mellersta raden i den inre 3x3{-}ytan (nollindexerad rad 2, positioner (2,1), (2,2), (2,3)) tomma (▢ ▢ ▢). Det finns inga X på den raden.} \\
\texttt{Fatta nu ett beslut. Vad tror du: är ditt rutnät och det andra rutnätet som beskrivs samma? Skriv "BESLUT: samma rutnät" om du tror att de är samma rutnät eller "BESLUT: olika rutnät" om du tror att de är olika rutnät. Lägg inte till något annat.} \\
            }
        }
    }
    & \\ \\

    \theutterance \stepcounter{utterance}  
    & \multicolumn{4}{p{0.6\linewidth}}{
        \cellcolor[rgb]{0.9,0.9,0.9}{
            \makecell[{{p{\linewidth}}}]{
                \texttt{\tiny{[P1$\rangle$GM]}}
                \texttt{BESLUT: olika rutnät} \\
            }
        }
    }
    & & \\ \\

    \theutterance \stepcounter{utterance}  
    & & & \multicolumn{2}{p{0.3\linewidth}}{
        \cellcolor[rgb]{0.9,0.9,0.9}{
            \makecell[{{p{\linewidth}}}]{
                \texttt{\tiny{[GM$|$GM]}}
                \texttt{continue} \\
            }
        }
    }
    & & \\ \\

    \theutterance \stepcounter{utterance}  
    & & & \multicolumn{2}{p{0.3\linewidth}}{
        \cellcolor[rgb]{0.9,0.9,0.9}{
            \makecell[{{p{\linewidth}}}]{
                \texttt{\tiny{[GM$|$GM]}}
                \texttt{success} \\
            }
        }
    }
    & & \\ \\

\end{supertabular}
}

\end{document}
